\documentclass[UTF8,zihao=-4]{ctexart}
\usepackage[a4paper,margin=2.2cm]{geometry}
\usepackage{array}
\usepackage{longtable}
\usepackage{hyperref}
\usepackage{verbatim}
\hypersetup{hidelinks}
\setlength{\parindent}{0pt}
\setlength{\parskip}{0.55em}
\setcounter{secnumdepth}{0}
\renewcommand{\arraystretch}{1.25}
\emergencystretch=4em
\sloppy
\title{03 体系结构与人机交互复习讲义}
\author{}
\date{}
\begin{document}
\maketitle
\tableofcontents
\newpage


\begin{quote}
范围：\texttt{03-01-软件设计基础与体系结构概念.pdf} 到 \texttt{03-04-人机交互设计.pdf}；结合 Exam/Review 中分层架构、MVC、物理包图、REST API 和 HCI 评价内容。
结构：先整理 PPT 主线，再放对应考试模板和快速背诵清单。
\end{quote}

\bigskip\hrule\bigskip

\section{3. 软件体系结构与人机交互}

\subsection{3.1 设计的任务}

设计回答“如何实现需求”。

与需求和编程的区别：

\begin{longtable}{|>{\raggedright\arraybackslash}p{0.24\textwidth}|>{\raggedright\arraybackslash}p{0.31\textwidth}|>{\raggedright\arraybackslash}p{0.31\textwidth}|}
\hline
\textbf{阶段} & \textbf{回答问题} & \textbf{关注模型} \\
\hline
\endfirsthead
\hline
\textbf{阶段} & \textbf{回答问题} & \textbf{关注模型} \\
\hline
\endhead
需求 & 做什么 & 用例图、概念类图 \\
\hline
设计 & 如何组织解决方案 & 包图、组件图、接口、协作 \\
\hline
编程 & 如何用代码实现 & 类、方法、语句、数据结构 \\
\hline
\end{longtable}

设计管理复杂度的基本手段：

\begin{itemize}
\item 关注点分离
\item 分解
\item 抽象
\item 模块化
\item 信息隐藏
\end{itemize}

\subsection{3.2 设计层次}

\subsubsection{高层设计：体系结构设计}

关注整个系统结构：

\begin{itemize}
\item 子系统
\item 分层
\item 组件
\item 部署节点
\item 连接方式
\item 技术栈约束
\end{itemize}

常用图：

\begin{itemize}
\item 组件图
\item 部署图
\item 高层包图
\end{itemize}

\subsubsection{中层设计}

关注包、模块、类之间如何协作。

常用图：

\begin{itemize}
\item 包图
\item 交互图
\item 设计类图
\end{itemize}

\subsubsection{低层设计}

关注单个类、方法、数据结构和算法。

内容：

\begin{itemize}
\item 属性
\item 方法
\item 参数
\item 返回值
\item 控制结构
\item 异常处理
\end{itemize}

\subsection{3.3 体系结构的组成}

体系结构由三类内容构成：

\begin{itemize}
\item 部件或组件
\item 连接件
\item 配置
\end{itemize}

解释：

\begin{itemize}
\item 组件承担计算、存储或业务职责。
\item 连接件定义组件交互方式，如方法调用、REST API、消息队列、数据库连接。
\item 配置说明组件和连接件如何组合成整体系统。
\end{itemize}

体系结构决策一旦落地，修改成本高。错误的层次、错误的依赖方向和错误的部署结构会形成长期技术债。

\subsection{3.4 4+1 视图}

4+1 视图用于从多个角度描述体系结构：

\begin{itemize}
\item 逻辑视图：类、包、领域对象、业务职责。
\item 开发视图：代码模块、工程结构、组件划分。
\item 进程视图：运行时进程、线程、并发、通信。
\item 物理视图：服务器、客户端、数据库、网络部署。
\item 场景视图：用关键用例串联并验证前四个视图。
\end{itemize}

答题时常用写法：

\begin{quote}
使用场景视图验证架构是否支持核心用例；使用逻辑视图说明业务分层；使用开发视图说明代码组织；使用进程视图说明运行时并发和通信；使用物理视图说明客户端、服务器和数据库部署。
\end{quote}

\subsection{3.5 体系结构风格}

体系结构风格定义：

\begin{itemize}
\item 组件词汇表
\item 连接件类型
\item 拓扑约束
\end{itemize}

\subsubsection{分层风格}

典型三层：

\begin{itemize}
\item 表示层
\item 业务逻辑层
\item 数据访问层
\end{itemize}

优点：

\begin{itemize}
\item 职责清晰。
\item 上层不用关心下层实现细节。
\item 便于替换界面、业务逻辑或数据存储。
\item 便于测试和维护。
\end{itemize}

缺点：

\begin{itemize}
\item 层间调用带来额外开销。
\item 简单功能也要穿过多层。
\item 层次划分错误会导致绕层访问和依赖混乱。
\end{itemize}

\subsubsection{MVC}

MVC 分离：

\begin{itemize}
\item Model：数据和业务状态。
\item View：界面显示。
\item Controller：接收输入并协调模型和视图。
\end{itemize}

Web 项目中常见对应关系：

\begin{itemize}
\item Controller：处理 HTTP 请求。
\item Service：执行业务逻辑。
\item Repository 或 DAO：访问数据。
\item Entity 或 PO：持久化对象。
\item DTO 或 VO：接口传输对象或界面展示对象。
\end{itemize}

\subsection{3.6 体系结构设计过程}

课程中的过程可概括为：

\begin{enumerate}
\item 分析需求。
\item 选择体系结构风格。
\item 设计逻辑结构。
\item 设计接口。
\item 设计物理部署。
\item 验证体系结构。
\item 评审并迭代。
\end{enumerate}

答题时不要只画图，要说明为什么选择这个结构。

示例：

\begin{quote}
企业信息系统强调可维护性、安全性和分工协作，适合采用分层架构。表示层处理用户交互，业务逻辑层封装业务规则，数据访问层封装持久化细节。上层依赖接口而非实现，以降低层间耦合。
\end{quote}

\subsection{3.7 接口与依赖倒置}

分层系统中要避免上层直接依赖下层实现类。

推荐结构：

\begin{itemize}
\item \texttt{service} 包声明业务接口。
\item \texttt{service.impl} 包提供业务实现。
\item \texttt{repository} 包声明数据访问接口。
\item \texttt{repository.impl} 包提供数据访问实现。
\item Controller 依赖 Service 接口。
\item Service 实现依赖 Repository 接口。
\end{itemize}

依赖倒置原则：

\begin{itemize}
\item 高层模块不依赖低层模块，二者都依赖抽象。
\item 抽象不依赖细节，细节依赖抽象。
\end{itemize}

考试代码题重点：

\begin{itemize}
\item 写出接口所在包名。
\item Controller 中通过构造器注入 Service。
\item ServiceImpl 中通过构造器注入 Repository。
\item Controller 不直接访问 Repository。
\end{itemize}

\subsection{3.8 体系结构验证}

验证关注：

\begin{itemize}
\item 核心用例是否能穿过架构完成。
\item 层间依赖方向是否正确。
\item 接口是否稳定。
\item 质量属性是否得到支撑。
\item 部署结构是否满足性能、安全和可靠性要求。
\item 是否存在架构坏味道。
\end{itemize}

常见架构坏味道：

\begin{itemize}
\item 上层绕过业务层直接访问数据库。
\item 所有业务逻辑堆在 Controller。
\item Service 只有转发，没有业务职责。
\item Repository 混入业务规则。
\item 包之间循环依赖。
\item 共享全局状态过多。
\end{itemize}

\bigskip\hrule\bigskip

\subsection{3.9 人机交互设计}

人机交互关注用户、任务、界面和系统反馈之间的关系。

可用性属性：

\begin{itemize}
\item 易学性
\item 效率
\item 易记性
\item 错误率与恢复能力
\item 满意度
\end{itemize}

易学性和效率存在取舍。新手需要清晰入口和提示，熟练用户需要快捷键、批量操作和命令式能力。

\subsection{3.10 人的认知特点}

设计原则：

\begin{itemize}
\item 识别优于回忆。
\item 图形信息在很多场景下比纯文本更易识别。
\item GUI 通过菜单、按钮和图标降低记忆负担。
\item CLI 对专家高效，但学习成本高。
\end{itemize}

\subsection{3.11 Fitts 定律}

Fitts 定律说明：

\begin{itemize}
\item 目标越近，操作越快。
\item 目标越大，操作越快。
\item 屏幕边缘和角落更容易命中。
\end{itemize}

设计推论：

\begin{itemize}
\item 高频按钮要放在易触达位置。
\item 危险按钮与常用按钮要保持距离。
\item 小图标按钮需要足够点击区域。
\end{itemize}

\subsection{3.12 用户类型}

\begin{itemize}
\item 新手：需要清晰导航、提示和容错。
\item 中级用户：需要稳定流程和效率提升。
\item 专家：需要快捷键、批处理和可配置操作。
\end{itemize}

好的界面同时服务不同熟练度用户：入口清楚，操作路径短，反馈及时，并提供快捷方式。

\subsection{3.13 直接操纵}

直接操纵的特点：

\begin{itemize}
\item 对象可见。
\item 操作可逆。
\item 反馈即时。
\item 用户有控制感。
\end{itemize}

优势：

\begin{itemize}
\item 学习时间短。
\item 错误易发现。
\item 用户满意度高。
\end{itemize}

需要处理的问题：

\begin{itemize}
\item 界面隐喻要符合用户心智模型。
\item 操作对象和系统状态要清晰。
\item 复杂系统需要良好导航结构。
\end{itemize}

\subsection{3.14 导航与反馈}

导航设计：

\begin{itemize}
\item 全局结构：用户知道自己在哪里、能去哪。
\item 局部结构：当前任务内的步骤清楚。
\item 按任务模型组织，而不是按技术模块组织。
\end{itemize}

反馈时间：

\begin{itemize}
\item 键盘、鼠标、光标反馈要非常快。
\item 简单频繁任务在约 1 秒内给出响应。
\item 更长操作要显示进度、状态或队列。
\end{itemize}

\subsection{3.15 黄金规则与 Nielsen 启发式}

三条黄金规则：

\begin{itemize}
\item 用户控制。
\item 减少记忆负担。
\item 保持一致。
\end{itemize}

Nielsen 十项启发式：

\begin{enumerate}
\item 系统状态可见。
\item 系统与现实世界匹配。
\item 用户控制和自由。
\item 一致性和标准。
\item 错误预防。
\item 识别优于回忆。
\item 灵活性和效率。
\item 美学和极简设计。
\item 帮助用户识别、诊断和恢复错误。
\item 帮助和文档。
\end{enumerate}

\subsection{3.16 HCI 设计过程与评估}

常用过程：

\begin{enumerate}
\item 理解用户和任务。
\item 建立界面概念模型。
\item 绘制线框图。
\item 制作低保真原型。
\item 用户测试。
\item 制作高保真原型。
\item 可用性评估。
\item 实现并迭代。
\end{enumerate}

评估方法：

\begin{itemize}
\item 启发式评估
\item 用户测试
\item Think-aloud
\item A/B 测试
\item 眼动分析
\end{itemize}

AI 交互设计补充：

\begin{itemize}
\item 显示 AI 的计划和依据。
\item 允许用户确认、撤销和覆盖。
\item 保留审计日志。
\item 在失败时给出可恢复路径。
\item 通过逐步授权建立信任。
\end{itemize}

\bigskip\hrule\bigskip

\bigskip\hrule\bigskip

\section{对应考试模板}

\subsection{7.1 分层接口代码模板}

题型：

\begin{quote}
给定某个情境，写出 Service 接口声明、Repository 方法声明、Controller 调用 Service 的代码片段，并体现依赖注入。
\end{quote}

作答要求：

\begin{itemize}
\item 写包名。
\item 写接口。
\item 写构造器注入。
\item Controller 调用 Service。
\item ServiceImpl 调用 Repository。
\item Controller 不直接调用 Repository。
\end{itemize}

示例：校园卡充值

\small
\begin{verbatim}
package edu.school.card.service;

import java.math.BigDecimal;
import java.util.List;

public interface RechargeService {
    RechargeResult recharge(RechargeCommand command);

    List<RechargeRecordView> findRechargeRecords(String cardNo);
}
\end{verbatim}
\normalsize

\small
\begin{verbatim}
package edu.school.card.repository;

import java.util.List;
import java.util.Optional;

public interface CampusCardRepository {
    Optional<CampusCard> findByCardNo(String cardNo);

    void save(CampusCard campusCard);
}
\end{verbatim}
\normalsize

\small
\begin{verbatim}
package edu.school.card.repository;

import java.util.List;

public interface RechargeRecordRepository {
    void save(RechargeRecord rechargeRecord);

    List<RechargeRecord> findByCardNo(String cardNo);
}
\end{verbatim}
\normalsize

\small
\begin{verbatim}
package edu.school.card.controller;

import edu.school.card.service.RechargeCommand;
import edu.school.card.service.RechargeResult;
import edu.school.card.service.RechargeService;
import edu.school.card.service.RechargeRecordView;
import java.util.List;
import org.springframework.web.bind.annotation.GetMapping;
import org.springframework.web.bind.annotation.PathVariable;
import org.springframework.web.bind.annotation.PostMapping;
import org.springframework.web.bind.annotation.RequestBody;
import org.springframework.web.bind.annotation.RequestMapping;
import org.springframework.web.bind.annotation.RestController;

@RestController
@RequestMapping("/api/recharges")
public class RechargeController {
    private final RechargeService rechargeService;

    public RechargeController(RechargeService rechargeService) {
        this.rechargeService = rechargeService;
    }

    @PostMapping
    public RechargeResult recharge(@RequestBody RechargeCommand command) {
        return rechargeService.recharge(command);
    }

    @GetMapping("/cards/{cardNo}")
    public List<RechargeRecordView> findRechargeRecords(
            @PathVariable String cardNo) {
        return rechargeService.findRechargeRecords(cardNo);
    }
}
\end{verbatim}
\normalsize

\small
\begin{verbatim}
package edu.school.card.service.impl;

import edu.school.card.repository.CampusCardRepository;
import edu.school.card.repository.RechargeRecordRepository;
import edu.school.card.service.RechargeCommand;
import edu.school.card.service.RechargeRecordView;
import edu.school.card.service.RechargeResult;
import edu.school.card.service.RechargeService;
import java.util.List;

public class RechargeServiceImpl implements RechargeService {
    private final CampusCardRepository campusCardRepository;
    private final RechargeRecordRepository rechargeRecordRepository;

    public RechargeServiceImpl(
            CampusCardRepository campusCardRepository,
            RechargeRecordRepository rechargeRecordRepository) {
        this.campusCardRepository = campusCardRepository;
        this.rechargeRecordRepository = rechargeRecordRepository;
    }

    @Override
    public RechargeResult recharge(RechargeCommand command) {
        CampusCard card = campusCardRepository.findByCardNo(command.cardNo())
                .orElseThrow(CardNotFoundException::new);
        card.recharge(command.amount());
        RechargeRecord record = RechargeRecord.success(
                command.cardNo(),
                command.amount());
        campusCardRepository.save(card);
        rechargeRecordRepository.save(record);
        return RechargeResult.success(record.id());
    }

    @Override
    public List<RechargeRecordView> findRechargeRecords(String cardNo) {
        return rechargeRecordRepository.findByCardNo(cardNo).stream()
                .map(RechargeRecordView::from)
                .toList();
    }
}
\end{verbatim}
\normalsize

答题说明：

\begin{itemize}
\item \texttt{controller} 是表示层入口。
\item \texttt{service} 是表示层与逻辑层之间的接口。
\item \texttt{repository} 是逻辑层与数据层之间的接口。
\item 构造器注入体现依赖注入。
\item 代码里出现的 DTO、Entity 和异常类可以只声明名称，重点是层间依赖方向。
\end{itemize}

\subsection{7.2 校园卡自助系统物理包图文字版}

题型：

\begin{quote}
画出校园卡自助系统物理包图。分为客户端和服务器端，体现分层概念，包含所有模块。三包和展示层在客户端，逻辑层和数据层在服务器端，使用 HTTP 通信和 REST API。
\end{quote}

考试作图要素：

\begin{itemize}
\item 客户端节点。
\item 服务器端节点。
\item 客户端包含 HTML、CSS、JavaScript 三包。
\item 展示层在客户端。
\item 逻辑层和数据层在服务器端。
\item 客户端与服务器通过 HTTP 通信。
\item 通信接口是 REST API。
\item 服务器端内部体现 Controller、Service、Repository、数据库。
\end{itemize}

文字图：

\small
\begin{verbatim}
<<device>> Client
  package html
  package css
  package javascript
  package presentation
    module LoginPage
    module RechargePage
    module RechargeRecordPage
    module LossReportPage
    module BalanceQueryPage
    module ProfilePage

  presentation --HTTP/REST API--> server.api

<<server>> Server
  package api
    module AuthController
    module RechargeController
    module BalanceController
    module LossReportController
    module CardController

  package service
    interface AuthService
    interface RechargeService
    interface BalanceService
    interface LossReportService
    interface CardService

  package service.impl
    class AuthServiceImpl
    class RechargeServiceImpl
    class BalanceServiceImpl
    class LossReportServiceImpl
    class CardServiceImpl

  package repository
    interface UserRepository
    interface CampusCardRepository
    interface RechargeRecordRepository
    interface LossReportRepository
    interface TransactionRepository

  package repository.impl
    class UserRepositoryImpl
    class CampusCardRepositoryImpl
    class RechargeRecordRepositoryImpl
    class LossReportRepositoryImpl
    class TransactionRepositoryImpl

  package domain
    class User
    class CampusCard
    class RechargeRecord
    class LossReport
    class Transaction

  package database
    table user
    table campus_card
    table recharge_record
    table loss_report
    table transaction
\end{verbatim}
\normalsize

依赖方向：

\small
\begin{verbatim}
presentation -> api -> service -> repository -> database
service.impl -> domain
repository.impl -> domain
\end{verbatim}
\normalsize

图上必须体现：

\small
\begin{verbatim}
Client.presentation -- HTTP + REST API --> Server.api
\end{verbatim}
\normalsize

\subsection{7.3 充值服务接口题模板}

题型：

\begin{quote}
针对充值服务，包括充值、查询充值记录，写出展示层和逻辑层之间接口，以及逻辑层和数据层之间接口。注意写出接口所在包名。
\end{quote}

展示层和逻辑层之间接口：

\small
\begin{verbatim}
package edu.school.card.service;

import java.util.List;

public interface RechargeService {
    RechargeResult recharge(RechargeCommand command);

    List<RechargeRecordView> findRechargeRecords(String cardNo);
}
\end{verbatim}
\normalsize

逻辑层和数据层之间接口：

\small
\begin{verbatim}
package edu.school.card.repository;

import java.util.Optional;

public interface CampusCardRepository {
    Optional<CampusCard> findByCardNo(String cardNo);

    void save(CampusCard campusCard);
}
\end{verbatim}
\normalsize

\small
\begin{verbatim}
package edu.school.card.repository;

import java.util.List;

public interface RechargeRecordRepository {
    void save(RechargeRecord rechargeRecord);

    List<RechargeRecord> findByCardNo(String cardNo);
}
\end{verbatim}
\normalsize

DTO 示例：

\small
\begin{verbatim}
package edu.school.card.service;

import java.math.BigDecimal;

public record RechargeCommand(String cardNo, BigDecimal amount, String payChannel) {
}
\end{verbatim}
\normalsize

\small
\begin{verbatim}
package edu.school.card.service;

public record RechargeResult(String recordId, String status, String message) {
}
\end{verbatim}
\normalsize

\small
\begin{verbatim}
package edu.school.card.service;

import java.math.BigDecimal;
import java.time.LocalDateTime;

public record RechargeRecordView(
        String recordId,
        String cardNo,
        BigDecimal amount,
        String status,
        LocalDateTime createdAt) {
}
\end{verbatim}
\normalsize

答题解释：

\begin{itemize}
\item \texttt{RechargeService} 是展示层调用逻辑层的接口。
\item \texttt{CampusCardRepository} 和 \texttt{RechargeRecordRepository} 是逻辑层调用数据层的接口。
\item \texttt{RechargeService} 不暴露数据库表结构。
\item \texttt{Repository} 不处理充值业务规则，只负责数据读取和保存。
\end{itemize}

\subsection{7.11 HCI 题模板}

题目要求评价界面：

\small
\begin{verbatim}
问题：充值按钮太小且远离金额输入区域。
原则：违反 Fitts 定律和效率原则。
影响：用户完成高频充值操作的时间增加，误触风险上升。
改进：增大按钮点击区域，把提交按钮放在金额输入区域附近，并对危险操作使用确认反馈。
\end{verbatim}
\normalsize

题目要求设计反馈：

\small
\begin{verbatim}
充值提交后，系统立即显示处理中状态；
支付成功后显示余额变化和充值记录编号；
支付失败时显示失败原因和重试入口；
长时间等待时显示进度或可取消操作。
\end{verbatim}
\normalsize

\bigskip\hrule\bigskip

\section{快速背诵清单}

\subsection{体系结构}

\begin{itemize}
\item 架构 = 组件 + 连接件 + 配置。
\item 4+1：逻辑、开发、进程、物理、场景。
\item 分层：表示层、业务逻辑层、数据访问层。
\item 依赖倒置：上层依赖接口，不依赖实现。
\end{itemize}

\subsection{HCI}

\begin{itemize}
\item 可用性：易学、效率、易记、错误恢复、满意度。
\item Fitts：目标越近越大，操作越快。
\item 识别优于回忆。
\item 黄金规则：用户控制、减少记忆负担、一致性。
\end{itemize}


\end{document}
